\documentclass[11pt]{article}
\usepackage[T1]{fontenc}
\usepackage[utf8]{inputenc}
\usepackage[a4paper,margin=27mm]{geometry}
\usepackage{lmodern,amsmath,amssymb,amsthm,microtype,booktabs,longtable,array}
\usepackage[authoryear,round]{natbib}
\usepackage{hyperref}
\hypersetup{hidelinks,pdftitle={Fixed cubic characters in natural-row moment estimates: After bramble\_trace},pdfauthor={}}
\setlength{\emergencystretch}{3em}
\setlength{\parskip}{3pt}
\newtheorem{theorem}{Theorem}[section]
\newtheorem{proposition}[theorem]{Proposition}
\newtheorem{lemma}[theorem]{Lemma}
\theoremstyle{remark}\newtheorem{remark}[theorem]{Remark}
\newcommand{\workbenchurl}{https://github.com/KokunoYumeto/zeta-and-l-functions-workbench}
\newcommand{\recordurl}{https://github.com/KokunoYumeto/mathematics-workbenches}
\newcommand{\N}{\mathrm N}
\DeclareMathOperator{\cont}{cont}
\DeclareMathOperator{\rad}{rad}
\providecommand{\tightlist}{\setlength{\itemsep}{0pt}\setlength{\parskip}{0pt}}
\title{Fixed cubic characters in natural-row moment estimates\\\large After bramble\_trace}
\author{}
\date{24 November 2026}
\expandafter\def\expandafter\UrlBreaks\expandafter{\UrlBreaks\do\a\do\b\do\c\do\d\do\e\do\f\do\g\do\h\do\i\do\j\do\k\do\l\do\m\do\n\do\o\do\p\do\q\do\r\do\s\do\t\do\u\do\v\do\w\do\x\do\y\do\z\do\A\do\B\do\C\do\D\do\E\do\F\do\G\do\H\do\I\do\J\do\K\do\L\do\M\do\N\do\O\do\P\do\Q\do\R\do\S\do\T\do\U\do\V\do\W\do\X\do\Y\do\Z\do\0\do\1\do\2\do\3\do\4\do\5\do\6\do\7\do\8\do\9}
\begin{document}
\maketitle
\begin{abstract}
bramble\_trace's inverse-product moments, natural-row density estimate, and smooth-amplifier bounds provide the starting estimates for this fixed-character applicability analysis. We treat the faithful cubic character of $\mathbf Q(\omega,\sqrt[3]{269})/\mathbf Q(\omega)$ and show that the principal fixed-twist convention does not exclude it: multiplication of the varying index absorbs the cubic twist with every zero mask included. Direct extraction and every fixed higher-power extraction give the absolute-counting exponent. We prove exact reconstruction from the different cube presentations, including all common divisors among inverse positions, and obtain the bound $Q^2+Y$ for the squared centrally normalized product. A shifted ordinary amplifier has a positive counting main term at the target, but its principal inducing rows contribute $P^{1/3}Q^2$ to its fourth moment and prevent an improvement in this extraction. We also prove the cost of an unrestricted coefficient block, give an explicit resonant coefficient, and explain why the stated density estimate permits a fixed exceptional inducing function. The current content cutoff and its remaining square-root cancellation problem are unchanged. The source analytic moments are explicitly granted; the complete elementary and algebraic deductions are proved here. The mathematics and text were produced by GPT-6 Astra in ChatGPT Pro mode and checked by separate instances of the same model; no human has checked this text.
\end{abstract}

\section{Main statements and exact scope}\label{sec:main}

A mean estimate over a varying family can be substantially stronger than absolute counting and still fail to give cancellation for one prescribed member. The distinction becomes especially concrete for cubic characters, because infinitely many natural indices induce the same primitive character. Those indices also carry different finite zero masks. This paper identifies exactly what the positive moment estimates of \citet{bt52v4,bt52v5,bt55v2} give when applied to one such character. It then supplies the deletion reconstruction and amplifier accounting needed to make two possible follow-up arguments valid and explicit.

The external analytic work is due to the cited sources. Its ingredients include the cubic large-sieve and Gauss-sum developments discussed by \citet{df26}, \citet{ddfs26}, and \citet{dr24}. We do not reprove those inputs or claim priority for a method. Our contribution here is the explicit application analysis, the complete conditional continuations, and the counterexamples to specified inferences. An upper-bound expression that does not improve an exponent is not an impossibility theorem for the underlying arithmetic problem.

\subsection{Arithmetic conventions and imported hypotheses}

Let $K=\mathbf Q(\omega)$, where $\omega^2+\omega+1=0$, and let $\mathfrak a=(269)$. All ideals are nonzero integral ideals. Fix a finite bad set $\mathcal S$ containing the primes above two and three and excluding $\mathfrak a$. All row and column ideals are prime to $\mathcal S$ unless an exception is stated. Fix a common column-deletion ideal $E$ divisible by $\mathfrak a$. We use the primary cubic-symbol convention and include every symbol's zero value at primes dividing its displayed modulus. Set $\psi=\chi_{\mathfrak a}$. Standard cubic reciprocity, Kummer theory, and unique factorization of ideals are background inputs.

Let $\mu_K$ be the ideal M\"obius function: it is zero on ideals divisible by a prime square and is $(-1)^r$ on a product of $r$ distinct prime ideals. For a fixed integer $k\ge1$, real heights $t_j$, smooth profiles $W_j$ supported in fixed compact annuli of $(0,\infty)$, and lengths $L_j\ge1$, put
\begin{align}
M_{h,j}(L_j;D)
 &=L_j^{-1/2}\sum_{(\mathfrak n,D)=1}
 \mu_K(\mathfrak n)\chi_h(\mathfrak n)(\N\mathfrak n)^{it_j}
 W_j(\N\mathfrak n/L_j),\label{eq:defM}\\
Z_h(\mathbf L;D)&=\prod_{j=1}^kM_{h,j}(L_j;D),\qquad
Y=\prod_{j=1}^kL_j,\qquad T=2+\max_j|t_j|.
\end{align}
Write $M_{\psi,j}=M_{\mathfrak a,j}$, $Z_\psi=Z_{\mathfrak a}$, and $A_\psi=Y^{1/2}Z_\psi$. Thus $A_\psi$ is the product before central normalization. The inverse product moment, granted in precisely the form used below, is
\begin{equation}
\sum_{0<\N h\le P}|Z_h(\mathbf L;D)|^2
\ll(PY\N D)^\epsilon T^J(P+P^{1/3}Y),
\qquad P\ge2,\quad L_j\ge1.
\tag{N}\label{eq:N}
\end{equation}
This is Corollary~3.10, (N1), of \citet{bt52v4,bt52v5}. It has no requirement $Y\le P$. The deletion, profiles, and heights are common to every row in each use. The number of factors is fixed before the loss and the finite profile-seminorm and height orders are selected. Constants have finite homogeneous dependence on the prescribed profile seminorms.

For a smooth annular profile $W_0$, define
\begin{equation}
S_h(Q;D)=Q^{-1/2}\sum_{(\mathfrak n,D)=1}
\chi_h(\mathfrak n)W_0(\N\mathfrak n/Q).
\end{equation}
We also grant the equal-factor case of the ordinary two-factor moment in \citet[Theorem~1 and Section~2]{bt55v2}:
\begin{equation}
\sum_{\substack{\N h\le P\\\chi_h\text{ induces a nonprincipal character}}}
 |S_h(Q;D)|^4\ll(PQ\N D)^\epsilon P,
 \qquad 1\le Q\le P^C.
\tag{O}\label{eq:O}
\end{equation}
Here $C$ is fixed, as are the common profiles and deletion. Fixed multiplicative changes of the row range are permitted. The source allows a common extra radical in a fixed polynomial range. In our amplifier application $D=E$ is fixed. Crucially, \eqref{eq:O} makes no assertion that includes principal inducing rows.

Write $N(\sigma,H;\chi_h)$ for the number of zeros of the primitive inducing $L$-function with real part at least $\sigma$ and imaginary part of absolute value at most $H$, counted with multiplicity. Finally, when discussing zeros we grant \citet[Theorem~5.7]{bt52v5} in its natural-row form: for every fixed $1/2<\sigma<1$ and every $\epsilon>0$, some fixed $\eta>0$ satisfies
\begin{equation}
\sum_{0<\N h\le P}N(\sigma,P^\eta;\chi_h)
\ll P^{D(\sigma)+\epsilon},\qquad
D(\sigma)=\frac{8-7\sigma}{3(2-\sigma)}.
\tag{D}\label{eq:D}
\end{equation}
Zeros belong to the primitive inducing functions, are counted with multiplicity, and are counted again for each natural presentation. Constants and $\eta$ are not claimed uniform at either endpoint of the open strip. The weighted-moment input for the unrestricted block is specified in Proposition~\ref{prop:block} and Section~\ref{sec:source}.

\subsection{The fixed character and the two continuation estimates}

\begin{theorem}[Fixed-character applicability and the limits of positive extraction]\label{thm:fixed}
The following statements hold with the preceding conventions.
\begin{enumerate}
\item The ideal $\mathfrak a$ is prime, $\N\mathfrak a=269^2=72361$, and $\psi$ is one faithful orientation of the Kummer character of $K(\sqrt[3]{269})/K$. The zero-extended identity
\[
\psi(\mathfrak n)\chi_h(\mathfrak n)=\chi_{\mathfrak a h}(\mathfrak n)
\]
embeds a fixed cubic twist in the natural family, by an injective fixed enlargement of the row range. The principal fixed-twist convention does not exclude this character.
\item Direct extraction of $h=\mathfrak a$ from \eqref{eq:N} gives
\[
|A_\psi(\mathbf L;E)|\ll_{\psi,E,k,\mathbf W,\epsilon}
Y^{1+\epsilon}T^J.
\]
Every fixed higher-power version of this extraction has the same exponent one. A larger row range does not improve it. The cube-free dyadic version, with $Y\le P$, does not permit an unbounded length while that fixed row remains in a fixed dyad.
\item The natural presentations with primitive inducing character $\psi$ have the form $h=\mathfrak a\mathfrak q^3$, with the fixed primary convention understood, and satisfy
\[
\chi_{\mathfrak a\mathfrak q^3}(\mathfrak n)
=\psi(\mathfrak n)\mathbf1_{(\mathfrak n,\mathfrak q)=1}.
\]
Their number below row norm $P$ is
\[
c_{\mathfrak a,\mathcal S}P^{1/3}+O_{\mathfrak a,\mathcal S}(P^{1/6}+1),
\qquad c_{\mathfrak a,\mathcal S}>0.
\]
Even idealized identical polynomial values on this many rows would yield only a squared normalized bound with positive term $Y$ from \eqref{eq:N}.
\item For every $b\in(1/2,1]$ there is a scalar family obeying \eqref{eq:N}, all its fixed-factor-count versions, and the stated cube-free $Y\le P$ moment, whose unnormalized fixed-row products have size proportional to $Y^b$. This is a countermodel only to inference from those positive inequalities: it is not a Möbius sequence or an Euler product and does not enforce the deletion identities.
\item A fixed zero of $L(s,\psi)$ at real part at least $\sigma$, once contained in the height window, contributes order $P^{1/3}$ to \eqref{eq:D}. Because
\[
D(\sigma)-\frac13=\frac{2(1-\sigma)}{2-\sigma}>0,
\]
this is compatible with that density upper bound for every fixed $\sigma<1$. No individual zero exclusion follows from this comparison.
\end{enumerate}
\end{theorem}

\begin{theorem}[Cube reconstruction and a shifted amplifier]\label{thm:continuation}
Assume \eqref{eq:N}. For $2\le Q\le\min_jL_j$,
\begin{align}
|Z_\psi(\mathbf L;E)|^2&\ll(QY)^\epsilon T^J(Q^2+Y),\label{eq:reconstruction}\\
|A_\psi(\mathbf L;E)|^2&\ll(QY)^\epsilon T^J(YQ^2+Y^2).
\end{align}
The proof uses the zero masks of the cube presentations and includes every common divisor among different inverse positions.

Let $\mathfrak b=\mathfrak a^2$, and suppose $W_0\ge0$ is not identically zero. There is $c_E>0$ such that
\begin{equation}
S_{\mathfrak b\mathfrak a}(Q;E)
=c_E\sqrt Q\int_0^\infty W_0(u)\,du+O_{E,W_0}(1).
\end{equation}
The inducing character of $\chi_{\mathfrak b h}$ is principal exactly for $h=\mathfrak a\mathfrak q^3$. Under \eqref{eq:O}, in its fixed polynomial range,
\begin{equation}
\sum_{\N h\le P}|S_{\mathfrak b h}(Q;E)|^4
\ll(PQ)^\epsilon(P+P^{1/3}Q^2).
\end{equation}
Under both \eqref{eq:N} and \eqref{eq:O}, for $P\ge\N\mathfrak a$, sufficiently large $Q$, and $Q\le P^C$ with fixed $C$,
\begin{equation}
|Z_\psi(\mathbf L;E)|^2
\ll(PQY)^\epsilon T^J
\frac{\sqrt{(P+P^{1/3}Y^2)(P+P^{1/3}Q^2)}}{Q}.
\label{eq:amplifier}
\end{equation}
Every inverse length is at least one; its fixed count is allowed to double in \eqref{eq:N}. The algebraic fraction on the right of \eqref{eq:amplifier} is at least $P^{1/3}Y$. This lower bound concerns the displayed upper-bound expression, not the value of the inverse sum.
\end{theorem}

\begin{proposition}[Unrestricted blocks and the scope of smooth amplification]\label{prop:block}
Fix an integer $j\ge2$. Suppose a weighted-moment upper bound has among its positive summands
\begin{equation}
P^{2/3-2/(3j)}U^{1/3+2/(3j)}L^{2/3-2/(3j)}.
\label{eq:weightedterm}
\end{equation}
This is the second summand of \citet[Proposition~7, (14.3)]{bt55v2}, whose hypotheses require fixed polynomial ranges in an ambient parameter, common profiles and zero masks, fixed divisor-bounded unrestricted coefficients, and the stated small-prime support for its smooth blocks. With unit-length ordinary factors, total smooth-amplifier length $S$, and unrestricted length $R$, one has $U=L=SR$, so \eqref{eq:weightedterm} equals $P^{2/3-2/(3j)}SR$.

If a nonzero centrally normalized amplifier has squared size at most $C_\delta S^{1+\delta}$, division by it leaves at least
\[
C_\delta^{-1}P^{2/3-2/(3j)}RS^{-\delta}
\]
in that numerical upper-bound expression. Under exact absolute counting $\ll S$, the resulting term is of order $P^{2/3-2/(3j)}R$; in fixed polynomial ranges the displayed $S^{-\delta}$ is an arbitrarily small-power qualification. Moreover the unrestricted coefficient class contains $c(\mathfrak n)=\overline{\psi(\mathfrak n)}$ on the fixed good support. For any nonzero nonnegative annular profile, its unnormalized fixed-$\psi$ sum has order $R$. This conclusion does not describe the Möbius subclass as resonant.
\end{proposition}

\begin{proposition}[Content weights and the transformed positive estimate]\label{prop:content}
The rational-content weight $\mu_K(\mathfrak n)\cont(\mathfrak n)^{it}$ is not multiplicative for generic real $t$, even on a product of distinct good prime ideals. Consequently it is not obtained by replacing the permitted norm heights in the source inverse product by content heights. The source's whole-tuple squarefree projection does not remove this defect. In a positive transformed moment with row scale $P=Y^2/H$ and outer scalar $H/Y$, the expression in \eqref{eq:N} becomes exactly
\[
\frac HY(P+P^{1/3}Y)=Y+(HY)^{2/3}.
\]
Adding the separate diagonal cost $H$ gives $H+Y+(HY)^{2/3}$, and for $H\ge\sqrt Y$ this is $O(H+(HY)^{2/3})$. This substitution supplies no strict improvement of that mixed exponent.
\end{proposition}

The significance of these statements is a precise division between applicability and strength. The desired character really belongs to the natural family. Exact deletion reconstruction is possible. A shifted amplifier with a positive counting main term is also possible. Their complete known costs nevertheless give the absolute-counting exponent in these extractions. No new zero-free boundary, smaller content cutoff, or square-root estimate for the remaining signed convolution is asserted.

\section{Proof overview}

The first step is arithmetic: cubic reciprocity identifies the character of the radicand 269 with the row of ideal norm 72361. Multiplication of the row index absorbs its fixed cubic twist. Positivity then isolates that row, but the second term in the natural moment contains a full factor of the total length. Repeating factors does not alter the exponent after taking the corresponding root. The presentation fiber has only order $P^{1/3}$ elements, and the natural density exponent remains strictly above $1/3$ throughout its stated open range.

The reconstruction argument addresses the deletion issue directly. In each Möbius position, every squarefree ideal has a unique factor supported on a given presentation ideal $\mathfrak q$. Expanding over these factors restores exactly the deleted terms. Different positions may use the same prime. After Cauchy--Schwarz, we fix the entire divisor tuple before applying the source moment. The row family then has common original profiles, heights and deletion, while its rescaled lengths stay at least one. Elementary ideal harmonic sums pay the remaining divisor costs.

For amplification, the fixed shift $\mathfrak b=\mathfrak a^2$ makes the ordinary factor principal at the target row and hence of size $\sqrt Q$. It also creates order $P^{1/3}$ principal inducing rows. Absolute counting bounds their combined fourth-moment cost by $P^{1/3}Q^2$. The nonprincipal fourth moment applies to the complement. Cauchy--Schwarz with the doubled inverse product gives the asserted formula, and multiplying its two second terms explains its length cost.

Finally, the unrestricted-block conclusion is a transparent exponent substitution, reinforced by a bounded coefficient that exactly cancels the target character. The content-weight example and the reverse-transform calculations separate additional coefficient and exponent issues from the already established moment estimates. All external analytic estimates remain stated assumptions of these deductions.

\section{Source scope and the exact division of credit}\label{sec:source}

The source statements have different families, length ranges and coefficient classes. The following inventory records the distinctions used in this paper. Suppressed factors in the last column are precisely the allowed small powers, polynomial height factors, and finite homogeneous profile seminorms; they are never an assertion of uniformity in a growing factor count.

\begin{center}\small
\begin{tabular}{@{}p{.24\textwidth}p{.42\textwidth}p{.25\textwidth}@{}}
\toprule
Source locator & Family and length range & Positive bound\\\midrule
Paper52 v4, Cor.~3.9, (I1) & Cube-free dyadic $\N h\asymp P$; $L_j\ge1$, $2\le Y\le P$ & $P+P^{1/3}Y$\\[4pt]
Paper52 v4, (I2) & Equal length $L$, fixed power $k$, $L^k\le P$ & $PL^k+P^{1/3}L^{2k}$ before normalization\\[4pt]
Paper52 v4/v5, Cor.~3.10, (N1) & All natural good rows; $P\ge2$, all $L_j\ge1$; no $Y\le P$ cap & $P+P^{1/3}Y$\\[4pt]
Paper52 v4/v5, (N2) & Nonprincipal inducing rows; inverse product $R$ times one full ordinary factor $X$; fixed polynomial ranges & $P+P^{2/3}R$\\[4pt]
Paper52 v5, Thm.~3.11, (M1) & Any fixed mixture of M\"obius and full ordinary factors; all natural rows; fixed polynomial ranges & $P+P^{1/3}Y$\\[4pt]
Paper55 v2, Thm.~1, (1.2) & Nonprincipal inducing rows; two full ordinary factors of arbitrary fixed polynomial lengths & $P$\\
\bottomrule
\end{tabular}
\end{center}

In (I2), comparison of its two terms gives the diagonal-size range $L\le P^{2/(3k)}$, including its endpoint. This is an upper-bound regime. The fixed finite twist in the Paper52 inverse and mixed statements is principal. Section~\ref{sec:character} explains why a fixed cubic twist can nonetheless be absorbed into its varying row parameter; that observation does not permit arbitrary fixed noncubic twists.

Theorem~2 of Paper52 v4/v5 is a varying-index statement with
\[
P^{12/17+\gamma}\le Y\le P^{2-\tau},\qquad
\N v,\N e,1+|t|\le P^\kappa,
\]
where $\gamma>0$, $0<\tau\le1/2$, and suitable positive $\kappa,\eta$ give the upper bound $P^{2/3-\eta}Y^{5/3}$. For principal fixed twist the lower exponent can be replaced by $1/2+\gamma$. The positive constants are not quantified there. The fixed row is still subject to an upper column-length scale. These assertions must not be converted into asymptotic one-character estimates by silently holding the row range fixed while its permitted lengths go to infinity.

The mixed statement (M1) permits coefficients $\mu_K$ or $1$, a common cubic orientation, and its stated projection of the whole tuple product onto squarefree ideals. It does not assert independent orientations at different positions, arbitrary coefficient blocks, prime-prefix restrictions, or independent row-dependent profiles. The source proof's common Mellin separation preserves the individual norm profiles; it does not automatically admit an unrelated arithmetic statistic.

Paper55's Theorem~4 extends the amplifier class to common coefficients of absolute value at most one supported on ideals all of whose prime divisors have norm at most a specified small power of the ambient parameter. In its moving-width notation the capacity is $n_1+n_2+3z\le M$, and its smoothness mesh depends on the requested loss. Proposition~7 includes one unrestricted divisor-bounded cofactor and charges it explicitly by four positive terms. Its numerical Corollary~4.1 assumes, writing its scalar parameter as $E_0$ to distinguish it from our deletion ideal,
\[
P=BE_0^3,\quad 1\le E_0\le B^{1/25},\quad
U\le B,\quad L\le B^{3/25},\quad R\le B^{1/20},
\]
and gives $E_0^{-2}\mathcal T\ll B^{2599/2400+\epsilon}$. The inverse product may range through $B^{3/25}$ only with the stated rough cofactor restriction. A general balanced inverse pair at the present target scale is not covered by deleting its complementary coefficient.

These estimates are used with their stated families, coefficient classes, common profiles, and deletion restrictions. The arguments below establish consequences of these hypotheses; they assume no strengthening of the underlying large-sieve, Poisson, or ordinary-moment inputs.

\section{The character family and direct extraction}\label{sec:character}

Since $269\equiv2\pmod3$, the rational prime 269 is inert in the Eisenstein field. Thus $\mathfrak a=(269)$ is prime of norm $269^2=72361$. Its primary generator is $-269$, which is congruent to one modulo nine. Cubic reciprocity on primary generators and the fact that $-1=(-1)^3$ give, away from the indicated bad primes,
\[
\chi_{\mathfrak a}(\mathfrak n)
=\left(\frac{\mathfrak n}{-269}\right)_3
=\left(\frac{-269}{\mathfrak n}\right)_3
=\left(\frac{269}{\mathfrak n}\right)_3.
\]
Kummer Frobenius identifies the last character with one faithful orientation of $K(\sqrt[3]{269})/K$. Complete conjugation, including profile and norm height, treats the other orientation in absolute-value estimates. Cubic multiplicativity, including the common-prime zeros, yields
\[
\psi(\mathfrak n)\chi_h(\mathfrak n)=\chi_{\mathfrak a h}(\mathfrak n).
\]
The map $h\mapsto\mathfrak a h$ is injective and multiplies its norm by 72361. Positivity therefore transfers a natural-row moment to the fixed cubic twist at this fixed enlarged range. The prime $\mathfrak a$ must remain outside the row bad set; it belongs in the common column deletion $E$. The rational radicand 269 must not be substituted for the row norm 72361. After cube factors are added, the index norm and primitive conductor norm also differ.

If a target convention initially includes the prime over two in the columns, an inverse factor can be split into its two squarefree sectors at $(2)$. The sector divisible by that prime is its exact Möbius factor times a coprime remaining sum at a fixed rescaled length. At fixed factor count, the resulting finitely many norm rescalings change constants, not asymptotic exponents. Bounded subunit lengths, if any, contain only finitely many terms and are handled directly. This finite restoration is distinct from changing the asymptotic column class.

Choose a fixed $P_0\ge\max(2,\N\mathfrak a)$. Retaining the target row in \eqref{eq:N} and restoring normalization gives
\[
|A_\psi|^2\ll
Y^{1+\epsilon}T^J(P_0+P_0^{1/3}Y)
\ll_{\psi,E}Y^{2+\epsilon}T^J.
\]
Taking a square root and renaming losses proves the second assertion of Theorem~\ref{thm:fixed}. For $k$ copies of one factor of length $L$, the same calculation gives $|A_\psi(L)|^{2k}\ll L^{2k+\epsilon}T^J$. The $2k$-th root again has exponent one. The theorem fixes $k$, its profile-seminorm orders and height exponent before $L$ varies, so no growing-order argument is included.

Both $P$ and $P^{1/3}Y$ are increasing in $P$. In power notation, $P=Y^u$ gives unnormalized exponent
\[
\max\left\{\frac{1+u}{2},1+\frac u6\right\}\ge1
\qquad(u\ge0).
\]
The cube-free dyadic family places the fixed cube-free row only in bounded dyads $P\asymp\N\mathfrak a$. Its $Y\le P$ condition then bounds the column length. Theorem~2's $Y\le P^{2-\tau}$ has the analogous limitation. The mixed estimate (N2) controls $U_\psi S_\psi$, and dividing by $S_\psi$ would require a lower bound it does not supply. A unit-length ordinary factor, when chosen nonzero, changes this only by a fixed constant and leaves exponent one. The ordinary amplifier in Section~\ref{sec:principal} constructs a genuine lower bound by a row shift and then pays its full moment cost.

\section{Presentation counts, a scalar countermodel, and density}\label{sec:fibers}

For every good $\mathfrak q$, the zero-extended symbol identity is
\begin{equation}
\chi_{\mathfrak a\mathfrak q^3}(\mathfrak m)
=\psi(\mathfrak m)\mathbf1_{(\mathfrak m,\mathfrak q)=1}.
\label{eq:frozen11}
\end{equation}
It fixes the inducing character but generally changes the polynomial. Kummer theory identifies equality of oriented characters with equality of radicands in $K^\times/K^{\times3}$. Ideal valuations then force the form $\mathfrak a\mathfrak q^3$; the fixed primary convention fixes the unit representative. The opposite orientation has core $\mathfrak a^2$, and has the same counting exponent.

The Eisenstein ring is a principal ideal domain with six units. Nonzero integral ideals are lattice elements modulo these units. A norm disk of parameter $x$ has area proportional to $x$ and boundary length $O(\sqrt x)$. Comparing its lattice cells with the disk gives the count with error $O(\sqrt x+1)$. Finite inclusion--exclusion at a fixed set of prime ideals multiplies its main coefficient by the positive product of $1-(\N\mathfrak p)^{-1}$. Consequently, with the fixed bad set absorbed into $F$,
\begin{equation}
\#\{\mathfrak n:\N\mathfrak n\le x,(\mathfrak n,F)=1\}
=c_Fx+O_F(x^{1/2}+1),\qquad c_F>0.
\label{eq:frozen12}
\end{equation}
For $x\ge1$, the extra one can be absorbed. Substitution of $x=(P/72361)^{1/3}$ proves the presentation-count formula in Theorem~\ref{thm:fixed}. Elementary summation by parts also gives
\begin{equation}
H_1(Q):=\sum_{\N\mathfrak d\le Q}(\N\mathfrak d)^{-1}\ll1+\log Q,
\qquad
H_2(Q):=\sum_{\N\mathfrak d\le Q}(\N\mathfrak d)^{-2}\ll1.
\label{eq:frozen13}
\end{equation}
Unique ideal factorization gives $\tau(\mathfrak q)^k\ll_{k,\epsilon}(\N\mathfrak q)^\epsilon$: bound each prime-power divisor factor by a chosen small power at sufficiently large prime norm and absorb the finitely many smaller primes in the constant.

Even if one artificially assumed identical normalized values on $cP^{1/3}$ presentations, dividing \eqref{eq:N} by this count would give only
\[
|Z(Y)|^2\ll(PY)^\epsilon T^J(P^{2/3}+Y).
\]
The positive $Y$ term gives the same unnormalized exponent. The different deletion masks require the further reconstruction in Section~\ref{sec:reconstruction}, not this artificial equality.

To prove the scalar countermodel, fix $b\in(1/2,1]$ and the bounded profile functional $\ell(W)=\int_0^\infty W(u)\,du/u$ on the fixed compact annulus. Define
\[
\mathcal M_h(L;W)=
\begin{cases}
L^{b-1/2}\ell(W),&h=\mathfrak a\mathfrak q^3,\\
0,&\text{otherwise}.
\end{cases}
\]
For any fixed product, the squared normalized value on the fiber is $Y^{2b-1}\prod_j|\ell(W_j)|^2$. Since $0<2b-1\le1$ and $Y\ge1$, its sum over the fiber is at most a profile constant times $P^{1/3}Y$, which is bounded by $P+P^{1/3}Y$. The functional's homogeneous bound supplies the allowed profile dependence. This argument works for every fixed factor count, including repetitions. In the cube-free family, the only nonzero row is $\mathfrak a$; whenever $Y\le P$, its squared normalized value is at most a profile constant times $Y\le P$. Yet its unnormalized product at the target is proportional to $Y^b$ for nonzero functionals. A square-root estimate with loss smaller than $b-1/2$ does not follow from these inequalities alone. The model does not enforce Möbius arithmetic, Euler products, or all deletion identities, so its conclusion is exactly the logical nonimplication stated in Theorem~\ref{thm:fixed}.

For the density assertion, let $\rho$ be a fixed zero of $L(s,\psi)$ with $\Re\rho\ge\sigma$. The positive height power in \eqref{eq:D} eventually contains $\rho$. Every natural presentation has that same primitive inducing function, independently of its finite polynomial mask, so the zero contributes the presentation count, with its fixed multiplicity. But direct subtraction gives
\[
\frac{8-7\sigma}{3(2-\sigma)}-\frac13
=\frac{2(1-\sigma)}{2-\sigma}>0.
\]
An order-$P^{1/3}$ contribution is therefore compatible with the stated upper bound. Letting $\sigma$ approach one does not alter this conclusion: the theorem supplies no uniform constants or common positive height exponent in that limit. The cube-free density theorem also cannot count the same fixed core in arbitrarily large dyads. This proves the remaining parts of Theorem~\ref{thm:fixed}.
\section{Reconstruction from cube presentations}\label{sec:reconstruction}

For \(Q\ge2\), let \(\mathcal Q(Q)\) consist of ideals \(\mathfrak q\) of norm at most \(Q\) prime to \(E\) and the fixed bad set. By \eqref{eq:frozen12},

\begin{equation}
|\mathcal Q(Q)|\gg_E Q.\label{eq:frozen21}
\end{equation}

For bounded \(Q\), decrease the implicit positive constant if necessary; the ideal \((1)\) is always present. Thus \eqref{eq:frozen21} holds uniformly for every \(Q\ge2\).

Fix one \(\mathfrak q\in\mathcal Q(Q)\). Every squarefree ideal \(\mathfrak n\) decomposes uniquely as

\[
\mathfrak n=\mathfrak d\mathfrak m,
\qquad \mathfrak d\mid\operatorname{rad}(\mathfrak q),
\qquad (\mathfrak m,\mathfrak q)=1,
\]

by putting in \(\mathfrak d\) exactly those prime factors of \(\mathfrak n\) that divide \(\mathfrak q\). In this decomposition \((\mathfrak d,\mathfrak m)=1\), so \(\mu_K(\mathfrak n)=\mu_K(\mathfrak d)\mu_K(\mathfrak m)\). All nonsquarefree terms have zero Möbius coefficient. Substituting this decomposition and using \eqref{eq:frozen11} yields the exact identity

\begin{equation}
M_{\psi,j}(L_j;E)
=\sum_{\mathfrak d\mid\operatorname{rad}(\mathfrak q)}
\mu_K(\mathfrak d)\psi(\mathfrak d)
(N\mathfrak d)^{-1/2+it_j}
M_{\mathfrak a\mathfrak q^3,j}(L_j/N\mathfrak d;E).
\label{eq:frozen22}
\end{equation}

The factor \((N\mathfrak d)^{-1/2}\) comes from the change in central normalization. The symbol on the right enforces \((\mathfrak m,\mathfrak q)=1\); it is not being replaced by \(\psi\) on that deleted set. In particular, \eqref{eq:frozen22} does not assert equality between differently deleted polynomial values.

Multiply \eqref{eq:frozen22} independently for all \(k\) inverse positions. For a tuple \(\mathbf d=(\mathfrak d_1,\ldots,\mathfrak d_k)\), put

\[
B(\mathbf d)=\prod_{j=1}^kN\mathfrak d_j,
\qquad \mathbf L/ N\mathbf d=(L_j/N\mathfrak d_j)_{j=1}^k.
\]

Then

\begin{equation}
Z_\psi(\mathbf L;E)
=\sum_{\mathfrak d_j\mid\operatorname{rad}(\mathfrak q),\ 1\le j\le k}
\left[\prod_{j=1}^k\mu_K(\mathfrak d_j)\psi(\mathfrak d_j)(N\mathfrak d_j)^{it_j}\right]
B(\mathbf d)^{-1/2}
Z_{\mathfrak a\mathfrak q^3}(\mathbf L/N\mathbf d;E).
\label{eq:frozen23}
\end{equation}

There is no condition \((\mathfrak d_i,\mathfrak d_j)=1\), and no condition between the summation ideals at two different inverse positions. A prime dividing \(\mathfrak q\) may occur in every \(\mathfrak d_j\). The only coprimality used is the within-position decomposition that is forced by that position's Möbius coefficient.

The number of summands in \eqref{eq:frozen23} is at most \(\tau(\mathfrak q)^k\). Since the coefficient in square brackets has modulus at most one, Cauchy--Schwarz gives

\begin{equation}
|Z_\psi|^2
\le\tau(\mathfrak q)^k
\sum_{\mathfrak d_j\mid\operatorname{rad}(\mathfrak q)}
B(\mathbf d)^{-1}
|Z_{\mathfrak a\mathfrak q^3}(\mathbf L/N\mathbf d;E)|^2.
\label{eq:frozen24}
\end{equation}

Sum \eqref{eq:frozen24} over \(\mathfrak q\in\mathcal Q(Q)\), use \eqref{eq:frozen21}, and enlarge the nonnegative divisor-tuple sum to all \(N\mathfrak d_j\le Q\). For a fixed tuple, the map \(\mathfrak q\mapsto\mathfrak a\mathfrak q^3\) is injective by unique ideal factorization. Its image lies in the natural-row range

\[
0<Nh\le P_Q:=N\mathfrak a\,Q^3.
\]

Consequently,

\begin{equation}
|Z_\psi|^2\ll_E Q^{-1+\eta}
\sum_{N\mathfrak d_j\le Q}B(\mathbf d)^{-1}
\sum_{0<Nh\le P_Q}
|Z_h(\mathbf L/N\mathbf d;E)|^2.
\label{eq:frozen25}
\end{equation}

This is the decisive place at which fixing the divisor tuple before applying \eqref{eq:N} keeps all profiles, heights, lengths, and deletions common to the rows. If \(Q\le\min_jL_j\), every resulting length satisfies

\begin{equation}
L_j/N\mathfrak d_j\ge L_j/Q\ge1.
\label{eq:frozen26}
\end{equation}

Thus every term in \eqref{eq:frozen25} is within the printed source length range. Its length product is \(Y/B(\mathbf d)\); no condition relating this product to \(P_Q\) is needed. All profiles remain exactly the original \(W_j\), all heights remain exactly \(t_j\), and the deletion remains exactly \(E\). In particular there is no row-dependent profile or enlarged deletion hidden in the moment application.

Apply \eqref{eq:N}, with exponent \(\eta\), to \eqref{eq:frozen25}. Since \(E\) and \(\mathfrak a\) are fixed,

\begin{equation}
|Z_\psi|^2
\ll Q^{-1+\eta}(Q^3Y)^\eta T^J
\left(Q^3\sum_{N\mathfrak d_j\le Q}B(\mathbf d)^{-1}
+QY\sum_{N\mathfrak d_j\le Q}B(\mathbf d)^{-2}\right).
\label{eq:frozen27}
\end{equation}

The two sums factor, respectively, as \(H_1(Q)^k\) and \(H_2(Q)^k\). Equation \eqref{eq:frozen13} therefore implies

\[
|Z_\psi|^2\ll Q^{4\eta}Y^\eta T^J
\{Q^2(1+\log Q)^k+Y\}.
\]

Choose \(\eta\) small enough in terms of the requested \(\epsilon\), and absorb the logarithm into \(Q^\epsilon\). This proves, uniformly for \(2\le Q\le\min_jL_j\),

\begin{equation}
|Z_\psi(\mathbf L;E)|^2
\ll_{\psi,E,k,\mathbf W,\epsilon}(QY)^\epsilon T^J(Q^2+Y).
\label{eq:frozenA}
\end{equation}

Since \(A_\psi=Y^{1/2}Z_\psi\), multiplication by \(Y\) gives the stated \(YQ^2+Y^2\) bound. This proves the reconstruction part of Theorem~\ref{thm:continuation}. It uses the cube rows with their zero masks. It neither uses prime ideals as the presentation family nor invokes the one-row shortcut. As a separate numerical sanity check, the one-row term of \eqref{eq:N} at fixed \(P\ge N\mathfrak a\) is stronger than \eqref{eq:frozenA}; that observation does not establish the reconstruction operation proved above.

\section{Principal rows in a shifted amplifier}\label{sec:principal}

Let \(\mathfrak b=\mathfrak a^2\). The primary cubic-symbol and Kummer convention fixed above implies that \(\chi_{\mathfrak b h}\) has principal inducing character exactly when every prime-ideal valuation of \(\mathfrak b h\) is divisible by three. At \(\mathfrak a\), this condition is \(v_{\mathfrak a}(h)+2\equiv0\pmod3\); at every other prime it is \(v_{\mathfrak p}(h)\equiv0\pmod3\). Since valuations are nonnegative, these conditions are equivalent to

\begin{equation}
h=\mathfrak a\mathfrak q^3.\label{eq:frozen31}
\end{equation}

This classifies the inducing character; the original symbol still vanishes at primes dividing \(\mathfrak a\mathfrak q\). The convention and Kummer equivalence are the stated reciprocity and Kummer background, not a newly established conductor formula.

For \(h=\mathfrak a\), the symbol in the amplifier is \(\chi_{\mathfrak a^3}\). Since \(E\) contains \(\mathfrak a\), it is one at every ideal that survives the common deletion. Smooth summation of \eqref{eq:frozen12} yields

\begin{equation}
S_{\mathfrak b\mathfrak a}(Q;E)
=Q^{-1/2}\sum_{(\mathfrak n,E)=1}W_0(N\mathfrak n/Q)
=c_E\sqrt Q\int_0^\infty W_0(u)\,du+O_{E,W_0}(1).
\label{eq:frozen32}
\end{equation}

For completeness, Stieltjes integration against \(c_Ex+O_E(\sqrt x)\), followed by integration by parts on the compact annular support, gives an unnormalized error \(O_{E,W_0}(\sqrt Q)\); central normalization gives the error in \eqref{eq:frozen32}. Nonnegativity and nonvanishing of \(W_0\) imply its integral is positive. Therefore, for all sufficiently large \(Q\),

\begin{equation}
|S_{\mathfrak b\mathfrak a}(Q;E)|^2\gg_{E,W_0}Q.
\label{eq:frozen33}
\end{equation}

For nonprincipal-inducing \(\chi_{\mathfrak b h}\), inject \(h\mapsto g=\mathfrak b h\) into the natural range \(Ng\le N\mathfrak b\,P\). Applying \eqref{eq:O} with the same \(E\) and \(W_0\), and using the permitted fixed enlargement of the range, gives

\begin{equation}
\sum_{\substack{Nh\le P\\\chi_{\mathfrak b h}\text{ nonprincipal inducing}}}
|S_{\mathfrak b h}(Q;E)|^4\ll(PQ)^\epsilon P.
\label{eq:frozen34}
\end{equation}

If the intended ordinary range is \(Q\le P^C\), it is also contained in the source range \(Q\le(N\mathfrak b P)^C\). All constants may depend on this fixed \(C\).

The remaining rows have \eqref{eq:frozen31}, so their number is \(O((P/N\mathfrak a)^{1/3})=O(P^{1/3})\), by elementary ideal counting. For every such row, including every additional deletion mask, absolute counting gives

\[
|S_{\mathfrak b h}(Q;E)|\le Q^{-1/2}\sum_{\mathfrak n}|W_0(N\mathfrak n/Q)|\ll_{W_0}\sqrt Q.
\]

Their total fourth-moment contribution is therefore \(O(P^{1/3}Q^2)\). Combining it with \eqref{eq:frozen34},

\begin{equation}
\sum_{Nh\le P}|S_{\mathfrak b h}(Q;E)|^4
\ll(PQ)^\epsilon(P+P^{1/3}Q^2).
\label{eq:frozen35}
\end{equation}

The extra term has been proved by counting principal rows; it has not been attributed to \eqref{eq:O}, whose statement excludes those rows.

\section{Amplified extraction and its numerical limitation}\label{amplified-extraction-and-its-numerical-limitation}

Assume \(P\ge N\mathfrak a\), so the row \(h=\mathfrak a\) occurs in the natural family. Nonnegativity, \eqref{eq:frozen33}, and Cauchy--Schwarz give

\begin{equation}
Q|Z_\psi(\mathbf L;E)|^2
\ll\sum_{Nh\le P}|Z_h(\mathbf L;E)|^2|S_{\mathfrak b h}(Q;E)|^2
\le\left(\sum_{Nh\le P}|Z_h|^4\right)^{1/2}
\left(\sum_{Nh\le P}|S_{\mathfrak b h}|^4\right)^{1/2}.
\label{eq:frozen41}
\end{equation}

To bound the first factor, use \eqref{eq:N} with \(2k\) inverse factors, repeating each of the original factors twice. Their product is \(Z_h^2\), so the modulus squared in \eqref{eq:N} is exactly \(|Z_h|^4\). The resulting length product is \(Y^2\); every individual length is still at least one. The profiles and heights are duplicated but remain common, and the deletion is still \(E\). No complex-conjugate factor class has been assumed. Thus

\begin{equation}
\sum_{Nh\le P}|Z_h|^4
\ll(PY^2N(E))^\eta T^{J_{2k}}(P+P^{1/3}Y^2).
\label{eq:frozen42}
\end{equation}

Insert \eqref{eq:frozen35} and \eqref{eq:frozen42} into \eqref{eq:frozen41}, decrease the source small exponents as necessary, and rename the harmless fixed height exponent. The result is

\begin{equation}
|Z_\psi(\mathbf L;E)|^2
\ll(PQY)^\epsilon T^J
\frac{\sqrt{(P+P^{1/3}Y^2)(P+P^{1/3}Q^2)}}{Q}.
\label{eq:frozenB}
\end{equation}

This holds only for sufficiently large \(Q\) to ensure \eqref{eq:frozen33}, for \(P\ge N\mathfrak a\), with all inverse lengths at least one, and within the printed polynomial ordinary length range of \eqref{eq:O}. Fixed counts, exponents, profiles, and range constants are chosen before asymptotic parameters.

Define the displayed algebraic factor by \(F(P,Q,Y)\). Since its factors are sums of nonnegative terms,

\begin{equation}
F(P,Q,Y)
\ge\frac{\sqrt{(P^{1/3}Y^2)(P^{1/3}Q^2)}}{Q}
=P^{1/3}Y.
\label{eq:frozen43}
\end{equation}

This is a lower bound on the numerical upper-bound expression produced by this extraction. It says that this particular deduction cannot make that expression smaller than \(P^{1/3}Y\), even by optimizing \(Q\). It is not a lower bound on \(|Z_\psi|^2\): an inequality \(X\le C F\), together with \(F\ge G\), does not imply \(X\ge cG\). Both moments may be wasteful, and the inverse polynomial may exhibit cancellation or even vanish. Multiplying \eqref{eq:frozenB} by \(Y\) gives an upper-bound expression for the unnormalized squared product with an algebraic floor \(P^{1/3}Y^2\). No new cubic cancellation estimate follows from this formula.

\section{The unrestricted-block cost and explicit resonance}\label{the-unrestricted-block-cost-and-explicit-resonance}

The unit-length ordinary profiles can be chosen with support near one and positive value at one, so that only the unit ideal survives and their target values are fixed nonzero constants. Thus they do not affect the length calculation.

For fixed \(j\ge2\), put

\[
\alpha=\frac23-\frac{2}{3j},\quad
\beta=\frac13+\frac{2}{3j},\quad
\gamma=\frac23-\frac{2}{3j}.
\]

The assumed positive summand is \(P^\alpha U^\beta L^\gamma\). The stated application has \(U=L=SR\), and \(\beta+\gamma=1\). Therefore the summand is exactly

\begin{equation}
P^\alpha(SR)^{\beta+\gamma}=P^\alpha SR.
\label{eq:frozen51}
\end{equation}

For a centrally normalized smooth amplifier of total length \(S\), maximal absolute counting gives squared size at most a constant times \(S\), or \(S^{1+\delta}\) when divisor losses are recorded as arbitrarily small powers. Dividing the positive upper-bound term \eqref{eq:frozen51} by the squared amplifier value can therefore leave no smaller algebraic term than a constant times \(P^\alpha R\). With the explicitly stated small-power convention, the precise version is

\begin{equation}
\frac{P^\alpha SR}{|\mathcal A_\psi|^2}
\ge C_\delta^{-1}P^\alpha R S^{-\delta}
\quad\text{if }0<|\mathcal A_\psi|^2\le C_\delta S^{1+\delta}.
\label{eq:frozen52}
\end{equation}

A zero amplifier permits no extraction at all. Formula \eqref{eq:frozen52} describes the first power of \(R\) at the level of powers used in the moment estimate; it does not silently discard the factor \(S^{-\delta}\). In ordinary fixed polynomial parameter ranges this factor is another arbitrarily small loss. With an absolute-counting bound \(\ll S\), no such qualification is necessary. In all cases this is a statement about the displayed positive upper-bound term, not a lower bound for every target sum.

There is also a direct obstruction to universal cancellation for an unrestricted coefficient block. Choose a smooth annular \(V\ge0\), not identically zero, and set

\[
c_R(\mathfrak n)=\overline{\psi(\mathfrak n)}
\mathbf1_{(\mathfrak n,E)=1}.
\]

These coefficients have absolute value at most one, so they lie in the unrestricted bounded coefficient class (and in any larger divisor-bounded class). On the surviving ideals \(|\psi(\mathfrak n)|=1\). Hence

\begin{equation}
\sum_{\mathfrak n}c_R(\mathfrak n)\psi(\mathfrak n)V(N\mathfrak n/R)
=\sum_{(\mathfrak n,E)=1}V(N\mathfrak n/R)
=c_E R\int_0^\infty V(u)\,du+O_{E,V}(\sqrt R).
\label{eq:frozen53}
\end{equation}

The unnormalized fixed-character sum is therefore \(\asymp R\), and the centrally normalized squared sum is \(\asymp R\). If the block includes a fixed norm height, absorb its inverse phase into \(c_R\); the modulus bound is unchanged. Alternatively, for a hard dyadic block use \(c_R(\mathfrak n)=\overline\psi(\mathfrak n)\mathbf1_{R<N\mathfrak n\le2R,(\mathfrak n,E)=1}\), giving the same order by \eqref{eq:frozen12}. This example concerns genuinely unrestricted coefficients; it makes no claim that Möbius coefficients or other arithmetic coefficient subclasses are resonant.

\section{The current content cutoff and its remaining target}\label{sec:content}

For an ideal $\mathfrak n$, its rational content $\cont(\mathfrak n)$ is the largest positive integer $r$ such that $\mathfrak n=(r)\mathfrak m$ with $\mathfrak m$ integral. In the Eisenstein principal ideal domain, the content of a generator $a+b\omega$ is $\gcd(a,b)$, unchanged by a unit. The companion joint-content analysis uses $m(\mathfrak d)=\mu_K(\mathfrak d)\psi(\mathfrak d)$, two balanced inverse positions, and a small ordinary cofactor. Its exact remaining sum can be stated without importing its reduction theorem here.

Choose fixed smooth cutoffs $b,\chi,\chi_d:[0,\infty)\to[0,1]$, equal to one on $[0,1]$ and zero on $[2,\infty)$. Fix $r>1$, $\kappa>4r$, and $W\in C_c^\infty((1,2))$ with
\[
\|(u\partial_u)^nW\|_\infty
\le C_WH_W^n n!\{\log(e+n)\}^{rn}\qquad(n\ge0)
\]
for fixed positive $C_W,H_W$. Put
\[
U=\sqrt{2X},\qquad L_X=\{\log(3X)\}^4
\{\log(e+\log(3X))\}^{\kappa},\qquad Y_c=X/\N\mathfrak c,
\]
and
\begin{align*}
\mathcal B_W(\mathfrak c,\mathfrak d_1,\mathfrak d_2;X)
={}&b(\N\mathfrak d_1/U)b(\N\mathfrak d_2/U)
W(\N\mathfrak c\N\mathfrak d_1\N\mathfrak d_2/X)\\
&\times\left[1-\chi_d(L_X\N\mathfrak d_1\N\mathfrak d_2/X)\right].
\end{align*}
Here $d_{2,K}(\mathfrak c)$ counts ordered ideal factorizations into two factors. With all ideals prime to the primes above three and 269, the remaining expression is
\begin{align}
R_{\gamma,W}(X)=-\frac12\sum_{\mathfrak c,\mathfrak d_1,\mathfrak d_2}
&\psi(\mathfrak c)d_{2,K}(\mathfrak c)\log\N\mathfrak c\,
m(\mathfrak d_1)m(\mathfrak d_2)
\mathcal B_W(\mathfrak c,\mathfrak d_1,\mathfrak d_2;X)\notag\\
&\times\chi\!\left(\frac{\cont(\mathfrak d_1)\cont(\mathfrak d_2)}{Y_c^\gamma}\right).
\label{eq:target}
\end{align}
The prime above two, present in this target convention, is restored by the finite squarefree-sector argument in Section~\ref{sec:character} when comparing its inverse factors to our source conventions. No coprimality between the two inverse ideals or their contents is imposed.

The theorem ``Specialization of the published zero-free input'' in the companion \emph{A smaller balanced prime-sum criterion for the cubic Riemann hypothesis: After Heath--Brown's 1982 identity} \citep{jointcontent} supplies the contextual parameter
\[
B_*=0.873345856106718576801184796065\ldots.
\]
Its strongest current unconditional specialization uses both zero-free exponents $\beta=\rho=B_*$ and has
\[
\gamma=\frac{2B_*-1}{3B_*}=0.284992750196395\ldots,
\qquad1-2\gamma=0.430014499607209\ldots.
\]
Under classical zeta RH, its corresponding specialization is
\[
\gamma=\frac{B_*-1/2}{2B_*-1/4}=0.249447399928862\ldots,
\qquad1-2\gamma=0.501105200142275\ldots.
\]
These specializations are contextual inputs from the cited companion theorem. They are not new consequences of Paper52 or Paper55. This paper neither reproves that zero-free input nor gives a stronger specialization. The remaining all-weight question is
\[
R_{\gamma,W}(X)=O_{W,\epsilon}(X^{1/2+\epsilon}).
\]

There are two distinct applicability issues. Even if the content cutoff were removed and common norm profiles separated, the arithmetic row would still be the single fixed $\mathfrak a$; the preceding extraction calculations already give exponent one. In addition, Mellin separation of the content cutoff inserts content heights rather than the allowed norm heights. These do not preserve the multiplicative coefficient identity in the source reverse transform.

For a complete explicit check, let $\pi=3+\omega$ and $\bar\pi=2-\omega$. Their norms are seven, the corresponding prime ideals are distinct and coprime, and
\[
(3+\omega)(2-\omega)=7,
\qquad\cont((\pi))=\cont((\bar\pi))=1,
\qquad\cont((7))=7.
\]
Each prime has Möbius value $-1$, while their product has value $1$. Therefore
\[
\mu_K((\pi))\cont((\pi))^{it}
\mu_K((\bar\pi))\cont((\bar\pi))^{it}=1,
\qquad
\mu_K((7))\cont((7))^{it}=7^{it}.
\]
For generic $t$ these differ. The product $(7)$ is already ideal-squarefree, so the source's whole-tuple squarefree projection does not cure the mismatch. Conversely, the target does not impose that the two inverse positions be coprime. The projections cannot be exchanged.

On squarefree ideal support, rational content may be recovered arithmetically from the rational norm, but that does not make it a uniformly bounded smooth annular norm profile. A new coefficient identity and a uniform seminorm estimate would be necessary. Rational cube presentations are no substitute: $\N(\mathfrak a(r)^3)=72361r^6$, and their character is still \eqref{eq:frozen11}, with only a further deletion. In particular, rational integers prime to the conductor do not create new primitive cubic characters. This proves the content assertion of Proposition~\ref{prop:content} while leaving open any genuinely new coefficient treatment.

\section{Exact transformed exponents}\label{sec:exponents}

With dual row scale $P=Y^2/H$, multiply the natural positive bound by $H/Y$. Direct algebra gives
\[
\frac HY\left(\frac{Y^2}H+
\left(\frac{Y^2}H\right)^{1/3}Y\right)
=Y+(HY)^{2/3}.
\]
Adding $H$ proves the remaining assertion of Proposition~\ref{prop:content}. If $H\ge\sqrt Y$, then $(HY)^{2/3}\ge Y$; hence the total is $O(H+(HY)^{2/3})$, with no strict saving relative to that expression.

The more detailed reverse-transform substitution is also exact. Set
\[
H=\frac{Y^2}{PGL},\qquad Z=\frac{Y}{GD},\qquad
C=VD(G/L)^2.
\]
Multiplying the source's four monomials by $PL$ gives
\begin{align*}
PLH&=Y^2G^{-1},\\
PL(HZ)^{2/3}&=P^{1/3}Y^2L^{1/3}G^{-4/3}D^{-2/3},\\
PLH^{1/3}Z^{5/6}&=P^{2/3}Y^{3/2}L^{2/3}G^{-7/6}D^{-5/6},\\
PLH^{1/3}Z^{7/9}C^{1/6}
&=P^{2/3}Y^{13/9}V^{1/6}L^{1/3}G^{-7/9}D^{-11/18}.
\end{align*}
Each identity follows by adding the exponents contributed by its displayed factors. These are the entries of the reverse-transform computation corresponding to Paper52 (R11). In particular the final term includes both $V^{1/6}$ and $D^{-11/18}$. The calculation is not an independent proof of the source analytic bounds to which the monomials belong.

For a second comparison, set $P=Y^u$, $1<u<2$. The unnormalized moment $PY+P^{1/3}Y^2$ has exponent
\[
Q(u)=\max\{u+1,2+u/3\}.
\]
Subtracting this exponent from $5/3+2u/3$ gives
\[
\frac53+\frac{2u}3-Q(u)
=\min\left\{\frac{2-u}3,\frac{u-1}3\right\}.
\]
On $4/3\le u\le3/2$, the second entry is the minimum, because $u-1\le2-u$. The saving is therefore exactly $(u-1)/3$. An application with a previously incurred loss of that size needs a further strict saving. An unspecified positive saving in a differently ranged theorem is not a numerical certificate that a specified loss has been surpassed.

For completeness, the finite-power coverage appearing in the inherited audit is elementary. The intervals
\[
I_q=[4/(3q),2/q],\qquad q=2,\ldots,32,
\]
have extreme endpoints $1$ and $1/24$. Consecutive intervals overlap because $2/(q+1)\ge4/(3q)$ is equivalent to $q\ge2$. Their union therefore covers $[1/24,1]$. This verifies that interval arithmetic, but does not confer a source theorem at variable or unbounded factor count.

\section{Method and scope}\label{sec:method}

The analytic hypotheses are the moments \eqref{eq:N} and \eqref{eq:O}, the density estimate \eqref{eq:D}, and the weighted-moment summand \eqref{eq:weightedterm}, with the source locators and restrictions stated above. The proofs establish the deletion identities, parameter costs, scalar countermodel, resonant coefficient, and content-weight obstruction. They do not reprove the external infinite analytic estimates.

The cube reconstruction, shifted-amplifier estimate, and unrestricted-block argument received a separate statement-only conditional reconstruction by the same model. The broader applicability deductions received targeted same-model checks. The accompanying \href{\workbenchurl}{\emph{Moment-scope exact checks} and \emph{Cube reconstruction exact checks}} verify rational exponent calculations and finite squarefree-decomposition identities. The proofs establish the continuous parameter claims; these computations check finite identities and transcription, not the analytic hypotheses.

The mathematics and text were produced by GPT-6 Astra in ChatGPT Pro mode and checked by separate instances of the same model. No human or proof-assistant verification is claimed. Original text and programs are dedicated to the public domain under \href{https://creativecommons.org/publicdomain/zero/1.0/}{CC0 1.0 Universal}; external publications retain their rights. The \href{\workbenchurl}{workbench repository} provides the accompanying material, and the \href{\recordurl}{project record repository} provides the project context.

These conditional deductions do not improve $B_*$, the current joint-content cutoff, or the square-root bound sought for \eqref{eq:target}. Further progress requires additional structure that survives localization to the fixed primitive character, a stronger localized estimate, or sufficient signed cancellation with an explicitly favorable loss calculation. None of those possibilities is ruled out here.

\begin{thebibliography}{99}
\bibitem[bramble\_trace(2026a)]{bt52v4}
bramble\_trace. (2026a, October 26). \emph{Short M\"obius product moments and explicit zero density for cubic Hecke characters} (Paper 52, Version 4) [Preprint]. rStation. \url{https://rstation.ai/problems/riemann_hypothesis/papers/52/v4}

\bibitem[bramble\_trace(2026b)]{bt52v5}
bramble\_trace. (2026b, October 26). \emph{Mixed product moments and natural zero density for cubic Hecke characters} (Paper 52, Version 5) [Preprint]. rStation. \url{https://rstation.ai/problems/riemann_hypothesis/papers/52/v5}

\bibitem[bramble\_trace(2026c)]{bt55v2}
bramble\_trace. (2026c). \emph{Cubic product moments with amplifiers on smooth ideals} (Paper 55, Version 2) [Preprint]. rStation. \url{https://rstation.ai/problems/riemann_hypothesis/papers/55/v2}

\bibitem[bramble\_trace(2026d)]{bt58v1}
bramble\_trace. (2026d). \emph{Polar subtraction and squarefree amplifiers in cubic Gauss means} (Paper 58, Version 1) [Preprint]. rStation. \url{https://rstation.ai/problems/riemann_hypothesis/papers/58/v1}

\bibitem[David et al.(2026)]{ddfs26}
David, C., de Faveri, A., Dunn, A., \& Stucky, J. (2026). \emph{Non-vanishing for cubic Hecke $L$-functions} (Version 3) [Preprint]. arXiv. \url{https://arxiv.org/abs/2410.03048v3}

\bibitem[de Faveri(2026)]{df26}
de Faveri, A. (2026). \emph{Optimal large sieve for fixed order characters} (Version 1) [Preprint]. arXiv. \url{https://arxiv.org/abs/2610.04045v1}

\bibitem[Dunn and Radziwi\l\l(2024)]{dr24}
Dunn, A., \& Radziwi\l\l, M. (2024). \emph{Bias in cubic Gauss sums: Patterson's conjecture} (Version 3) [Preprint]. arXiv. \url{https://arxiv.org/abs/2109.07463v3}
\bibitem[Smaller balanced prime-sum criterion(2026)]{jointcontent}
\emph{A smaller balanced prime-sum criterion for the cubic Riemann hypothesis: After Heath--Brown's 1982 identity}. (2026, November 24). [Companion workbench paper]. \href{\workbenchurl}{Workbench repository}. Theorem ``Specialization of the published zero-free input.''

\end{thebibliography}
\end{document}
